\begin{table}[!htbp]
\centering
\caption{Sensitivity of PACE Eligibility Effects to Additional Classification Errors}
\label{tab:pace-classification-stress}
\small
\begin{tabular}{lrrr}
\toprule
\makecell{Additional error budget\\(share of all students)} & Students & \makecell{Attendance effect\\range (pp)} & \makecell{95\% confidence\\envelope (pp)} \\
\midrule
0.0\% & 0 & [-6.84, -2.68] & [-7.49, -2.02] \\[4pt]
0.1\% & 176 & [-7.86, -1.62] & [-8.68, -0.80] \\[4pt]
0.5\% & 880 & [-11.59, 2.41] & [-12.65, 3.47] \\[4pt]
1.0\% & 1,759 & [-15.91, 7.20] & [-17.16, 8.46] \\[4pt]
\bottomrule
\end{tabular}
\par\vspace{0.2cm}
\begin{minipage}{\textwidth}\footnotesize
Notes: The sample and specification are those in Table~\ref{tab:pace-operational}. Each row allows all unresolved eligibility classifications to vary and up to the indicated number of additional, currently determinate, binary outcomes to be wrong. The errors are allocated to students to maximize or minimize the fitted IV coefficient. The confidence envelope also allows the robust standard error to increase. These budgets are sensitivity assumptions, not estimated error rates or extrapolations from the selected applicant validation sample.
\end{minipage}
\end{table}
